\HLASourcePage{20}{19}
\thispagestyle{hla-body}
\begin{hlafixedpage}
  \HLAFixedBlock{90}{69.79}{432}{\rmfamily}{9.96}{12}{%
    The arguments to the services described in this document will use different views of the entities depending\linebreak
    on a particular RTI implementation. For clarity, the first 11 clauses of this document refer to only a generic\linebreak
    view, known as a ``designator,'' when referring to these entities, and the API annexes refer to the actual\linebreak
    implementation-specific arguments. Clause 12 ties the two concepts together.%
  }

  \HLAFixedBlock{90}{129.78}{432}{\rmfamily}{9.96}{12}{%
    Some of the arguments to the services described in this document are in the form of a group that represents\linebreak
    some association. These groupings can take several forms (depending on the requirements of the respective\linebreak
    service), each of which is described below. The groupings may be implemented slightly differently in each\linebreak
    API, but all implementations shall support the same semantics. These groupings are pairs, sets, collections,\linebreak
    constrained sets of pairs, and collections of pairs. Collections are sets that permit duplication of elements.\linebreak
    They are used in cases in which duplication is unlikely and not really an error. As such, using a collection\linebreak
    rather than a set retains the same effects, but without the expensive checks to eliminate duplicates that using\linebreak
    a set would require. In these cases, it is preferred not to penalize the majority of users with the less efficient\linebreak
    set just to catch errors that are anticipated to be infrequent.%
  }

  \HLAFixedBlock{90}{249.78}{432}{\rmfamily}{9.96}{12}{%
    The following data are needed for the implementation of a running RTI and federation executions:%
  }

  \HLAFixedDashList{90}{267.78}{432}{0}{%
    \item FOM Document Data (FDD): FOM information (e.g., class, attribute, parameter names) used by the\linebreak
      RTI at runtime. Each federation execution needs FDD. FDD is provided to the RTI starting with the\linebreak
      Management Object Model (MOM) and Initialization Module (MIM) and FOM modules specified\linebreak
      in the \emph{Create Federation Execution} service invocation. FDD will also be provided by the MIM that\linebreak
      is automatically provided by the RTI or optionally provided in the \emph{Create Federation Execution}\linebreak
      service invocation. Additional FDD may be provided as joining federates specify additional FOM\linebreak
      modules.
  }

  \HLAFixedBlock{90}{363.78}{432}{\rmfamily}{9.96}{12}{%
    The following data may be needed for the implementation of a running RTI and federation executions:%
  }

  \HLAFixedDashList{90}{381.78}{432}{0}{%
    \item RTI initialization data (RID): RTI vendor-specific information needed to run an RTI. If required,\linebreak
      RID are supplied when an RTI is initialized.
  }

  \HLAFixedBlock{90}{417.78}{432}{\rmfamily}{9.96}{12}{%
    For all joined federate-initiated services in this specification (except 4.2, 4.5, 4.6, and 4.11), an implied\linebreak
    supplied argument is a joined federate's conjointment with a federation execution. For all RTI-initiated\linebreak
    services, an implied supplied argument is also a joined federate's connection to a federation execution. The\linebreak
    manner in which these arguments are actually provided to the services is dependent on the RTI\linebreak
    implementation and the particular API programming language; therefore, it is not shown in the service\linebreak
    descriptions. Also, for the RTI-initiated services, some implicit preconditions are not stated explicitly\linebreak
    because the RTI is assumed to be well behaved. For many RTI-initiated services, the ``Federate is joined''\linebreak
    precondition presumes the ``Federate is connected'' precondition.%
  }

  \HLAFixedBlock{90}{525.78}{432}{\rmfamily}{9.96}{12}{%
    As mentioned earlier, all RTI-initiated services are denoted with a \textdagger{} (printer's dagger) after the service\linebreak
    name.%
  }

  \HLAFixedBlock{90}{561.78}{432}{\rmfamily}{9.96}{12}{%
    If an RTI-invoked service is not the result of a joined federate-invoked service, but the RTI-invoked service\linebreak
    has a user-supplied tag as a mandatory argument, the user-supplied tag shall be present in the service\linebreak
    invocation but empty. For example, in the Java API, the tag shall be an empty (zero-length) array. Examples\linebreak
    of RTI-invoked services with empty user-supplied tags are Provide Attribute Value Update (via Auto\linebreak
    Provide switch), Request Attribute Ownership (via Unconditional Divest resign directive), and MOM\linebreak
    HLAmanager.HLAfederate.HLAreport interactions.%
  }

  \HLAAnchor{clause:1.6}
  \pdfbookmark[1]{1.6 Organization of this document}{bookmark:clause:1.6}
  \HLAFixedBlock{90}{651.38}{250}{\sffamily\bfseries\raggedright}{10.98}{13}{%
    1.6 Organization of this document%
  }

  \HLAFixedBlock{90}{676.81}{432}{\rmfamily}{9.96}{12}{%
    The seven HLA service groups are specified in Clause 4 through Clause 10. Each service is described using\linebreak
    several components:%
  }

  \HLAFixedDashList{90}{706.81}{432}{0}{%
    \item Name and description: service name and narrative describing the functionality of the service.
  }
\end{hlafixedpage}
\clearpage
